\chapter{Coupled Fields}
\section{General}
The \xel{pdeList} element can be followed by a \xel{couplingList} element,
which can be used to supply information on the coupling between different
PDEs. Within this block we distinguish between iterative coupling
(\xel{iterative}) and direct coupling (\xel{direct}).
\section{Iterative Coupling}
 For iterative coupling,
the method can be \xatv{RHS} (for right hand side coupling) and \xatv{Matrix}
for
a matrix coupling algorithm.
Then, for each involved pair of coupling PDEs, we have
to specify for each PDE the coupling parameters. 
The coupling is specified by the coupling input \xat{quantity}, the
interface \xat{type} - which can be \xatv{node}(Coupling),
\xatv{surface}(Coupling) or \xatv{region}(Coupling) -  and the \xat{name}
of
the interface, which has to correspond with the according interface
\xat{type}. At the end, we have to define
the stopping criteria and the maximum number of iterations
within the section \xel{nonlLinear}. Therefore, the general structure 
looks as follows


\begin{Verbatim}[commandchars=@\[\]]
@PYaZ[<couplingList]@PYaZ[>]
  @PYaZ[<iterative]@PYaZ[>]
    @PYaZ[<Pde1Pde2] @PYaQ[method=]@PYaB["RHS"]@PYaZ[>]
      @PYaZ[<Pde1]@PYaZ[>]
        @PYaZ[<coupling] @PYaQ[quantity=]@PYaB["..."] @PYaQ[type=]@PYaB["..."] @PYaQ[name=]@PYaB["..."]@PYaZ[/>]
      @PYaZ[</Pde1>]
      @PYaZ[<Pde2]@PYaZ[>]
        @PYaZ[<coupling] @PYaQ[quantity=]@PYaB["..."] @PYaQ[type=]@PYaB["..."] @PYaQ[name=]@PYaB["..."]@PYaZ[/>]
        @PYaZ[<coupling] @PYaQ[quantity=]@PYaB["..."] @PYaQ[type=]@PYaB["..."] @PYaQ[name=]@PYaB["..."]@PYaZ[/>]
      @PYaZ[</Pde2>]
    @PYaZ[</Pde1Pd2>]

    @PYaZ[<Pde3Pde4] @PYaQ[method=]@PYaB["Matrix"]@PYaZ[>]
      @PYaZ[<Pde3]@PYaZ[>]
        @PYaZ[<coupling] @PYaQ[quantity=]@PYaB["..."] @PYaQ[type=]@PYaB["..."] @PYaQ[name=]@PYaB["..."]@PYaZ[/>]
        @PYaZ[<coupling] @PYaQ[quantity=]@PYaB["..."] @PYaQ[type=]@PYaB["..."] @PYaQ[name=]@PYaB["..."]@PYaZ[/>]
      @PYaZ[</Pde3>]
      @PYaZ[<Pde4]@PYaZ[>]
        @PYaZ[<coupling] @PYaQ[quantity=]@PYaB["..."] @PYaQ[type=]@PYaB["..."] @PYaQ[name=]@PYaB["..."]@PYaZ[/>]
      @PYaZ[</Pde4>]
    @PYaZ[</Pde4Pde4>]

    @PYaZ[<nonLinear] @PYaQ[logging=]@PYaB["no"]@PYaZ[>]
      @PYaZ[<stopCrit] @PYaQ[value=]@PYaB["1e-3"] @PYaQ[quantity=]@PYaB["..."] @PYaQ[l2Norm=]@PYaB["..."]@PYaZ[/>]
      @PYaZ[<stopCrit] @PYaQ[value=]@PYaB["1e-2"] @PYaQ[quantity=]@PYaB["..."] @PYaQ[l2Norm=]@PYaB["..."]@PYaZ[/>]
      @PYaZ[<stopCrit] @PYaQ[value=]@PYaB["1e-5"] @PYaQ[quantity=]@PYaB["..."] @PYaQ[l2Norm=]@PYaB["..."]@PYaZ[/>]
      @PYaZ[<maxNumIters]@PYaZ[>] 20 @PYaZ[</maxNumIters>]
    @PYaZ[</nonLinear>]
  @PYaZ[</iterative>]
@PYaZ[</couplingList>]
\end{Verbatim}


\noindent
For the stopping criteria, we have currently the following names:

\begin{Verbatim}[commandchars=@\[\]]
    @PYaZ[<stopCrit] @PYaQ[value=]@PYaB["1e-3"] @PYaQ[quantity=]@PYaB["mechDisplacement"] @PYaQ[l2Norm=]@PYaB["rel"]@PYaZ[/>]
    @PYaZ[<stopCrit] @PYaQ[value=]@PYaB["1e-3"] @PYaQ[quantity=]@PYaB["smoothDisplacement"] @PYaQ[l2Norm=]@PYaB["rel"]@PYaZ[/>]
    @PYaZ[<stopCrit] @PYaQ[value=]@PYaB["1e-3"] @PYaQ[quantity=]@PYaB["elecForceVWP"] @PYaQ[l2Norm=]@PYaB["rel"]@PYaZ[/>]
    @PYaZ[<stopCrit] @PYaQ[value=]@PYaB["1e-3"] @PYaQ[quantity=]@PYaB["magForceLorentz"] @PYaQ[l2Norm=]@PYaB["rel"]@PYaZ[/>]
    @PYaZ[<stopCrit] @PYaQ[value=]@PYaB["1e-2"] @PYaQ[quantity=]@PYaB["mechForce"] @PYaQ[l2Norm=]@PYaB["rel"]@PYaZ[/>]
    @PYaZ[<stopCrit] @PYaQ[value=]@PYaB["1e-2"] @PYaQ[quantity=]@PYaB["acouForce"] @PYaQ[l2Norm=]@PYaB["rel"]@PYaZ[/>]
\end{Verbatim}

Instead of using a relative accuracy $\epsilon^{rel}$, i.e. $\|u_{k+1}-u_k\|_2 <
\epsilon^{rel} \|u_{k+1}\|_2$ you can as well
truncate the iterations by choosing an absolute accuracy $\epsilon^{abs}$, i.e.
$\|u_{k+1}-u_k\|_2 < \epsilon^{abs}$ setting {\sf
l2Norm="abs"}. Note that the specification of the stopping criterias
is optional: The coupling is considered as converged if either
{\xel{maxNumIters} } iterations have been performed or if {\bf all} stopping
criterias were met.

\newpage
\subsection{Specifications within the single fields}
Currently we have definitions for electrostatics, mechanics, magnetics and
acoustics.

\begin{table}[htb]
\begin{center}
\caption{Electrostatics}
\bigskip
\begin{tabular}{ll}
\toprule
input coupling & type \\
\midrule
nodeCoupling / regionCoupling & mechDisplacement \\
nodeCoupling / regionCoupling & smoothDisplacement \\
\bottomrule
\end{tabular}
\end{center}
\end{table}

\begin{table}[htb]
\begin{center}
\caption{Mechanics}
\bigskip
\begin{tabular}{ll}
\toprule
input coupling    & type \\
\midrule
nodeCoupling / regionCoupling     & elecForceVWP \\
nodeCoupling / regionCoupling    & magForceVWP \\
regionCoupling    & magForceLorentz \\
interfaceCoupling & acouForce \\
\bottomrule
\end{tabular}
\end{center}
\end{table}

\begin{table}[htb]
\begin{center}
\caption{Magnetics}
\bigskip
\begin{tabular}{ll}
\toprule
input coupling    & type \\
\midrule
nodeCoupling / regionCoupling & mechDisplacement \\
nodeCoupling / regionCoupling & smoothDisplacement \\
\bottomrule
\end{tabular}
\end{center}
\end{table}

\begin{table}[htb]
\begin{center}
\caption{Acoustics}
\bigskip
\begin{tabular}{ll}
\toprule
input coupling    & type \\
\midrule
interfaceCoupling      & mechForce \\
\bottomrule
\end{tabular}
\end{center}
\end{table}

\newpage
\subsection{ElecMech}
The two involved fields are the electrostatic and the mechanical field.
Often we need a additional PDE for the smoothing of the mesh (moving/deforming 
body problem).

\bigskip
\noindent
{\it Example:}

\begin{Verbatim}[commandchars=@\[\]]
@PYaZ[<couplingList]@PYaZ[>]
  @PYaZ[<iterative]@PYaZ[>]
    @PYaZ[<elecMech] @PYaQ[method=]@PYaB["RHS"]@PYaZ[>]
      @PYaZ[<electrostatic]@PYaZ[>]
        @PYaZ[<coupling] @PYaQ[quantity=]@PYaB["smoothDisplacement"] @PYaQ[type=]@PYaB["region"] @PYaQ[name=]@PYaB["air1"]@PYaZ[/>]
        @PYaZ[<coupling] @PYaQ[quantity=]@PYaB["smoothDisplacement"] @PYaQ[type=]@PYaB["region"] @PYaQ[name=]@PYaB["air2"]@PYaZ[/>]
      @PYaZ[</electrostatic>]
      @PYaZ[<mechanic]@PYaZ[>]
        @PYaZ[<coupling] @PYaQ[quantity=]@PYaB["elecForceVWP"] @PYaQ[type=]@PYaB["node"] @PYaQ[name=]@PYaB["electrode1"]@PYaZ[/>]
      @PYaZ[</mechanic>]
      @PYaZ[<smooth]@PYaZ[>]
        @PYaZ[<coupling] @PYaQ[quantity=]@PYaB["mechDisplacement"] @PYaQ[type=]@PYaB["node"] @PYaQ[name=]@PYaB["electrode1"]@PYaZ[/>]
      @PYaZ[</smooth>]
    @PYaZ[</elecMech>]
    @PYaZ[<nonLinear] @PYaQ[logging=]@PYaB["no"]@PYaZ[>]
      @PYaZ[<stopCrit] @PYaQ[value=]@PYaB["1e-2"] @PYaQ[quantity=]@PYaB["mechDisplacement"] @PYaQ[l2Norm=]@PYaB["rel"]@PYaZ[/>]
      @PYaZ[<stopCrit] @PYaQ[value=]@PYaB["1e-3"] @PYaQ[quantity=]@PYaB["elecForceVWP"] @PYaQ[l2Norm=]@PYaB["abs"]@PYaZ[/>]
      @PYaZ[<maxNumIters]@PYaZ[>] 20 @PYaZ[</maxNumIters>]
    @PYaZ[</nonLinear>]
  @PYaZ[</iterative>]>
@PYaZ[</couplingList>]
\end{Verbatim}


\subsection{MagMech}
The two involved fields are the magnetic and the mechanical field.
Often we need a additional PDE for the smoothing of the mesh (moving/deforming 
body problem).

\bigskip
\noindent
{\it Example:}

\begin{Verbatim}[commandchars=@\[\]]
@PYaZ[<couplingList]@PYaZ[>]
  @PYaZ[<iterative]@PYaZ[>]
    @PYaZ[<magMech] @PYaQ[method=]@PYaB["RHS"]@PYaZ[>]
      @PYaZ[<magnetic]@PYaZ[>]
        @PYaZ[<coupling] @PYaQ[quantity=]@PYaB["smoothDisplacement"] @PYaQ[type=]@PYaB["region"] @PYaQ[name=]@PYaB["air1"]@PYaZ[/>]
        @PYaZ[<coupling] @PYaQ[quantity=]@PYaB["smoothDisplacement"] @PYaQ[type=]@PYaB["region"] @PYaQ[name=]@PYaB["air2"]@PYaZ[/>]
      @PYaZ[</magnetic>]
      @PYaZ[<mechanic]@PYaZ[>]
        @PYaZ[<coupling] @PYaQ[quantity=]@PYaB["magForceLorentz"] @PYaQ[type=]@PYaB["region"] @PYaQ[name=]@PYaB["plate"]@PYaZ[/>]
      @PYaZ[</mechanic>]
      @PYaZ[<smooth]@PYaZ[>]
        @PYaZ[<coupling] @PYaQ[quantity=]@PYaB["mechDisplacement"] @PYaQ[type=]@PYaB["node"] @PYaQ[name=]@PYaB["surf1"]@PYaZ[/>]
      @PYaZ[</smooth>]
      @PYaZ[<nonLinear] @PYaQ[logging=]@PYaB["no"]@PYaZ[>]
        @PYaZ[<stopCrit] @PYaQ[value=]@PYaB["1e-2"] @PYaQ[quantity=]@PYaB["mechDisplacement"] @PYaQ[l2Norm=]@PYaB["rel"]@PYaZ[/>]
        @PYaZ[<stopCrit] @PYaQ[value=]@PYaB["1e-3"] @PYaQ[quantity=]@PYaB["magForceLorentz"] @PYaQ[l2Norm=]@PYaB["abs"]@PYaZ[/>]
      @PYaZ[<maxNumIters]@PYaZ[>] 20 @PYaZ[</maxNumIters>]
    @PYaZ[</nonLinear>]
  @PYaZ[</iterative>]
@PYaZ[</couplingList>]
\end{Verbatim}



\subsection{AcouMech}
The two involved fields are the Acoustics and the mechanical field.

\bigskip
\noindent
{\it Example:}

\begin{Verbatim}[commandchars=@\[\]]
@PYaZ[<couplingList]@PYaZ[>]
  @PYaZ[<iterative]@PYaZ[>]
    @PYaZ[<acouMech] @PYaQ[method=]@PYaB["RHS"]@PYaZ[>]
      @PYaZ[<acoustic]@PYaZ[>]
        @PYaZ[<coupling] @PYaQ[quantity=]@PYaB["mechForce"] @PYaQ[type=]@PYaB["interface"] @PYaQ[name=]@PYaB["surf1"]@PYaZ[/>]
      @PYaZ[</acoustic>]
      @PYaZ[<mechanic]@PYaZ[>]
        @PYaZ[<coupling] @PYaQ[quantity=]@PYaB["acouForce"] @PYaQ[type=]@PYaB["interface"] @PYaQ[name=]@PYaB["surf1"]@PYaZ[/>]
      @PYaZ[</mechanic>]
    @PYaZ[</acouMech>]
    @PYaZ[<nonLinear] @PYaQ[logging=]@PYaB["no"]@PYaZ[>]
      @PYaZ[<stopCrit] @PYaQ[value=]@PYaB["1e-2"] @PYaQ[quantity=]@PYaB["mechDisplacement"] @PYaQ[l2Norm=]@PYaB["rel"]@PYaZ[/>]
      @PYaZ[<maxNumIters]@PYaZ[>] 20 @PYaZ[</maxNumIters>]
    @PYaZ[</nonLinear>]
  @PYaZ[</AcouMech>]
@PYaZ[</couplingList>]
\end{Verbatim}

\section{Direct Coupling}
Here we differ between an indirect and a direct coupling. Currently we have
just implemented the iterative coupling. Note that the coupling
between these two field can only be done via \xml{interface} coupling,
as the element normals are needed for calculating the coupling forces.

\bigskip
\noindent
{\it Example:}

\begin{Verbatim}[commandchars=@\[\]]
@PYaZ[<couplingList]@PYaZ[>]
  @PYaZ[<direct]@PYaZ[>]
    @PYaZ[<coupledpde]@PYaZ[>]
      @PYaZ[<regionList]@PYaZ[>]
        @PYaZ[<region] @PYaQ[name=]@PYaB["..."]@PYaZ[/>]
      @PYaZ[</regionList>]
      @PYaZ[<storeResults]@PYaZ[>]
        @PYaZ[<elemResult]@PYaZ[>]
          @PYaZ[<elemList]@PYaZ[>]
            @PYaZ[<elems] @PYaQ[name=]@PYaB["..."]@PYaZ[/>]
          @PYaZ[</elemList>]
        @PYaZ[</elemResult>]
        @PYaZ[<surfElemResult]@PYaZ[>]
          @PYaZ[<surfElemList]@PYaZ[>]
            @PYaZ[<surfElems] @PYaQ[name=]@PYaB["..."]@PYaZ[/>]
          @PYaZ[</surfElemList>]
        @PYaZ[</surfElemResult>]
      @PYaZ[</storeResults>]
      @PYaZ[<materialDataType] @PYaQ[type=]@PYaB["..."]@PYaZ[/>]
      @PYaZ[<nonLinear]@PYaZ[>]
        @PYaZ[<lineSearch]@PYaZ[/>]
        @PYaZ[<incStopCrit]@PYaZ[>] ... @PYaZ[</incStopCrit>]
        @PYaZ[<resStopCrit]@PYaZ[>] ... @PYaZ[</resStopCrit>]
        @PYaZ[<maxNumIters]@PYaZ[>] ... @PYaZ[</maxNumIters>]
      @PYaZ[</nonLinear>]
    @PYaZ[</coupledpde>]
  @PYaZ[</direct>]
@PYaZ[</couplingList>]
\end{Verbatim}


\subsection{PiezoDirect}
The two involved fields are the Electrostatic and the mechanical field.

\bigskip
\noindent
{\it Example:}

\begin{Verbatim}[commandchars=@\[\]]
@PYaZ[<couplingList]@PYaZ[>]
  @PYaZ[<direct]@PYaZ[>]
    @PYaZ[<piezoDirect]@PYaZ[>]
      @PYaZ[<regionList]@PYaZ[>]
        @PYaZ[<region] @PYaQ[name=]@PYaB["..."]@PYaZ[/>]
        @PYaZ[<region] @PYaQ[name=]@PYaB["..."]@PYaZ[/>]
      @PYaZ[</regionList>]
      @PYaZ[<storeResults]@PYaZ[>]
        @PYaZ[<elemResult] @PYaQ[type=]@PYaB["mechStress"]@PYaZ[>]
          @PYaZ[<allRegions]@PYaZ[/>]
        @PYaZ[</elemResult>]
      @PYaZ[</storeResults>]
    @PYaZ[</piezoDirect>]
  @PYaZ[</direct>]
@PYaZ[</couplingList>]
\end{Verbatim}


\subsection{AcouMechDirect}
The two involved fields are the Acoustics and Mechanical Field.

\bigskip
\noindent
{\it Example:}

\begin{Verbatim}[commandchars=@\[\]]
@PYaZ[<couplingList]@PYaZ[>]
  @PYaZ[<direct]@PYaZ[>]
    @PYaZ[<acouMechDirect]@PYaZ[>]
      @PYaZ[<surfRegionList]@PYaZ[>]
        @PYaZ[<surfRegion] @PYaQ[name=]@PYaB["..."]@PYaZ[/>]
      @PYaZ[</surfRegionList>]
      @PYaZ[<ncInterfaceList]@PYaZ[>]
        @PYaZ[<ncInterface] @PYaQ[name=]@PYaB["..."]@PYaZ[/>]
        @PYaZ[<ncInterface] @PYaQ[name=]@PYaB["..."]@PYaZ[/>]
      @PYaZ[</ncInterfaceList>]
    @PYaZ[</acouMechDirect>]
  @PYaZ[</direct>]
@PYaZ[</couplingList>]
\end{Verbatim}


\subsection{ThermoElectricDirect}
The two involved fields are the heatConduction and Electrostatic Field.

\bigskip
\noindent
{\it Example:}

\begin{Verbatim}[commandchars=@\[\]]
@PYaZ[<couplingList]@PYaZ[>]
  @PYaZ[<direct]@PYaZ[>]
    @PYaZ[<thermoElectricDirect]@PYaZ[>]
      @PYaZ[<regionList]@PYaZ[>]
        @PYaZ[<region] @PYaQ[name=]@PYaB["..."]@PYaZ[/>]
      @PYaZ[</regionList>]
    @PYaZ[</thermoElectricDirect>]
  @PYaZ[</direct>]
@PYaZ[</couplingList>]
\end{Verbatim}


\subsection{ThermoMechDirect}
The two involved fields are the heatConduction and Mechanical Field.

\bigskip
\noindent
{\it Example:}

\begin{Verbatim}[commandchars=@\[\]]
@PYaZ[<couplingList]@PYaZ[>]
  @PYaZ[<direct]@PYaZ[>]
    @PYaZ[<thermoMechDirect]@PYaZ[>]
      @PYaZ[<regionList]@PYaZ[>]
        @PYaZ[<region] @PYaQ[name=]@PYaB["..."]@PYaZ[/>]
      @PYaZ[</regionList>]
    @PYaZ[</thermoMechDirect>]
  @PYaZ[</direct>]
@PYaZ[</couplingList>]
\end{Verbatim}

